% GENERATED FILE. Edit canonical lessons, then run npm run generate:books.
% canonical-source-hash: fnv1a64:62994451d6c58f40
% canonical-lessons: TE-C304-prakrti, TE-C304-krida, TE-C304-atagadu, TE-C304-kriket, TE-C304-cadarangam, TE-R304-first-pass-words-from-chapters-296-299, TE-R304-second-pass-words-from-chapters-300-304

\chapter{Nature, Sport, Player, Cricket, Chess}
\label{ch:te-nature-sport-player-cricket-chess}
\hlchaptermodality{304}

\begin{chapteropening}
\textbf{By the end of this chapter:} \emph{I can say nature, a sport, a player, cricket and chess.}

Complete the last lesson of chapter 304: I can say nature, a sport, a player, cricket and chess.
\end{chapteropening}

\section[\texorpdfstring{\te{ప్రకృతి} (prakṛti)}{ప్రకృతి (prakṛti)}]{\te{ప్రకృతి} (prakṛti) — nature}

\label{lesson:TE-C304-prakrti}



\begin{quote}
Before the new one: say the Telugu for sugarcane, then the Telugu for worship.
\end{quote}

\subsection*{You'll want to know: \te{ప్రకృతి}}
\textbf{\te{ప్రకృతి}} — \emph{prakṛti} — \textquotedblleft{}nature\textquotedblright{}.

\begin{grammarlens}[title={What you've built}]
One more word for this chapter.
\end{grammarlens}

\subsection*{Your turn}
\begin{itemize}
\raggedright
  \item \emph{Say it:} \emph{prakṛti}
  \item \emph{Say it:} \emph{prakṛti}, once more
  \item \emph{Say it:} say \emph{prakṛti}
  \item \emph{Recall:} say the Telugu for sugarcane, then the Telugu for worship, then say \emph{prakṛti} again
\end{itemize}

\begin{tcolorbox}[breakable,colback=teal!4,colframe=teal!35!black,title={Before you move on}]
What does \te{ప్రకృతి} mean? (\textquotedblleft{}Nature\textquotedblright{}.) Say it once more.
\end{tcolorbox}
\section[\texorpdfstring{\te{క్రీడ} (krīḍa)}{క్రీడ (krīḍa)}]{\te{క్రీడ} (krīḍa) — a sport}

\label{lesson:TE-C304-krida}



\begin{quote}
Before the new one: say the Telugu for worship, then the Telugu for nature.
\end{quote}

\subsection*{You'll want to know: \te{క్రీడ}}
\textbf{\te{క్రీడ}} — \emph{krīḍa} — \textquotedblleft{}a sport\textquotedblright{}.

\begin{grammarlens}[title={What you've built}]
One more word for this chapter.
\end{grammarlens}

\subsection*{Your turn}
\begin{itemize}
\raggedright
  \item \emph{Say it:} \emph{krīḍa}
  \item \emph{Say it:} \emph{krīḍa}, once more
  \item \emph{Say it:} say \emph{krīḍa}
  \item \emph{Recall:} say the Telugu for worship, then the Telugu for nature, then say \emph{krīḍa} again
\end{itemize}

\begin{tcolorbox}[breakable,colback=teal!4,colframe=teal!35!black,title={Before you move on}]
What does \te{క్రీడ} mean? (\textquotedblleft{}A sport\textquotedblright{}.) Say it once more.
\end{tcolorbox}
\section[\texorpdfstring{\te{ఆటగాడు} (āṭagāḍu)}{ఆటగాడు (āṭagāḍu)}]{\te{ఆటగాడు} (āṭagāḍu) — a player}

\label{lesson:TE-C304-atagadu}



\begin{quote}
Before the new one: say the Telugu for nature, then the Telugu for a sport.
\end{quote}

\subsection*{You'll want to know: \te{ఆటగాడు}}
\textbf{\te{ఆటగాడు}} — \emph{āṭagāḍu} — \textquotedblleft{}a player\textquotedblright{}.

\begin{grammarlens}[title={What you've built}]
One more word for this chapter.
\end{grammarlens}

\subsection*{Your turn}
\begin{itemize}
\raggedright
  \item \emph{Say it:} \emph{āṭagāḍu}
  \item \emph{Say it:} \emph{āṭagāḍu}, once more
  \item \emph{Say it:} say \emph{āṭagāḍu}
  \item \emph{Recall:} say the Telugu for nature, then the Telugu for a sport, then say \emph{āṭagāḍu} again
\end{itemize}

\begin{tcolorbox}[breakable,colback=teal!4,colframe=teal!35!black,title={Before you move on}]
What does \te{ఆటగాడు} mean? (\textquotedblleft{}A player\textquotedblright{}.) Say it once more.
\end{tcolorbox}
\section[\texorpdfstring{\te{క్రికెట్} (krikeṭ)}{క్రికెట్ (krikeṭ)}]{\te{క్రికెట్} (krikeṭ) — cricket}

\label{lesson:TE-C304-kriket}



\begin{quote}
Before the new one: say the Telugu for a sport, then the Telugu for a player.
\end{quote}

\subsection*{You'll want to know: \te{క్రికెట్}}
\textbf{\te{క్రికెట్}} — \emph{krikeṭ} — \textquotedblleft{}cricket\textquotedblright{}.

\begin{grammarlens}[title={What you've built}]
One more word for this chapter.
\end{grammarlens}

\subsection*{Your turn}
\begin{itemize}
\raggedright
  \item \emph{Say it:} \emph{krikeṭ}
  \item \emph{Say it:} \emph{krikeṭ}, once more
  \item \emph{Say it:} say \emph{krikeṭ}
  \item \emph{Recall:} say the Telugu for a sport, then the Telugu for a player, then say \emph{krikeṭ} again
\end{itemize}

\begin{tcolorbox}[breakable,colback=teal!4,colframe=teal!35!black,title={Before you move on}]
What does \te{క్రికెట్} mean? (\textquotedblleft{}Cricket\textquotedblright{}.) Say it once more.
\end{tcolorbox}
\section[\texorpdfstring{\te{చదరంగం} (cadaraṅgaṁ)}{చదరంగం (cadaraṅgaṁ)}]{\te{చదరంగం} (cadaraṅgaṁ) — chess}

\label{lesson:TE-C304-cadarangam}



\begin{quote}
Before the new one: say the Telugu for a player, then the Telugu for cricket.
\end{quote}

\subsection*{You'll want to know: \te{చదరంగం}}
\textbf{\te{చదరంగం}} — \emph{cadaraṅgaṁ} — \textquotedblleft{}chess\textquotedblright{}.

\begin{grammarlens}[title={What you've built}]
One more word for this chapter.
\end{grammarlens}

\subsection*{Your turn}
\begin{itemize}
\raggedright
  \item \emph{Say it:} \emph{cadaraṅgaṁ}
  \item \emph{Say it:} \emph{cadaraṅgaṁ}, once more
  \item \emph{Say it:} say \emph{cadaraṅgaṁ}
  \item \emph{Recall:} say the Telugu for a player, then the Telugu for cricket, then say \emph{cadaraṅgaṁ} again
\end{itemize}

\begin{tcolorbox}[breakable,colback=teal!4,colframe=teal!35!black,title={Before you move on}]
What does \te{చదరంగం} mean? (\textquotedblleft{}Chess\textquotedblright{}.) Say it once more.
\end{tcolorbox}
\section[(dialogue)]{Words from chapters 296-299 — a first pass}

\label{lesson:TE-R304-first-pass-words-from-chapters-296-299}



\begin{quote}
Say the words for cricket and chess.
\end{quote}

\subsection*{Your turn}
Say these aloud:

\begin{itemize}
\raggedright
  \item hard, special, convenient, useful, local
  \item stubborn, stingy, active, naughty, honesty
  \item blunt, crooked, flat, hollow, dim
  \item chickpeas, mustard seeds, coriander leaves, curry leaves, buttermilk
\end{itemize}

\begin{tcolorbox}[breakable,colback=teal!4,colframe=teal!35!black,title={Before you move on}]
Hard? (\textbf{\te{కఠినమైన}}, \emph{kaṭhinamaina}.) Special? (\textbf{\te{ప్రత్యేక}}, \emph{pratyēka}.) Convenient? (\textbf{\te{అనుకూలమైన}}, \emph{anukūlamaina}.) Useful? (\textbf{\te{ఉపయోగకరమైన}}, \emph{upayōgakaramaina}.) Local? (\textbf{\te{స్థానిక}}, \emph{sthānika}.) Stubborn? (\textbf{\te{మొండి}}, \emph{moṇḍi}.) Stingy? (\textbf{\te{పిసినారి}}, \emph{pisināri}.) Active? (\textbf{\te{చురుకైన}}, \emph{curukaina}.) Naughty? (\textbf{\te{అల్లరి}}, \emph{allari}.) Honesty? (\textbf{\te{నిజాయితీ}}, \emph{nijāyitī}.) Blunt? (\textbf{\te{మొద్దు}}, \emph{moddu}.) Crooked? (\textbf{\te{వంకర}}, \emph{vaṅkara}.) Flat? (\textbf{\te{చదునైన}}, \emph{cadunaina}.) Hollow? (\textbf{\te{బోలు}}, \emph{bōlu}.) Dim? (\textbf{\te{మసక}}, \emph{masaka}.) Chickpeas? (\textbf{\te{శనగలు}}, \emph{śanagalu}.) Mustard seeds? (\textbf{\te{ఆవాలు}}, \emph{āvālu}.) Coriander leaves? (\textbf{\te{కొత్తిమీర}}, \emph{kottimīra}.) Curry leaves? (\textbf{\te{కరివేపాకు}}, \emph{karivēpāku}.) Buttermilk? (\textbf{\te{మజ్జిగ}}, \emph{majjiga}.)
\end{tcolorbox}
\section[(dialogue)]{Words from chapters 300-304 — a second pass}

\label{lesson:TE-R304-second-pass-words-from-chapters-300-304}



\begin{quote}
Say the words for shoes and a chain.
\end{quote}

\subsection*{Your turn}
Say these aloud:

\begin{itemize}
\raggedright
  \item shoes, a chain, earrings, a button, a handkerchief
  \item a dictionary, a cart, an auto-rickshaw, a taxi, petrol
  \item an ox, a calf, a sheep, a wolf, a crocodile
  \item a thorn, a creeper, bamboo, sugarcane, worship
  \item nature, a sport, a player, cricket, chess
\end{itemize}

\begin{tcolorbox}[breakable,colback=teal!4,colframe=teal!35!black,title={Before you move on}]
Shoes? (\textbf{\te{బూట్లు}}, \emph{būṭlu}.) A chain? (\textbf{\te{గొలుసు}}, \emph{golusu}.) Earrings? (\textbf{\te{చెవిపోగులు}}, \emph{cevipōgulu}.) A button? (\textbf{\te{గుండీ}}, \emph{guṇḍī}.) A handkerchief? (\textbf{\te{చేతిరుమాలు}}, \emph{cētirumālu}.) A dictionary? (\textbf{\te{నిఘంటువు}}, \emph{nighaṇṭuvu}.) A cart? (\textbf{\te{బండి}}, \emph{baṇḍi}.) An auto-rickshaw? (\textbf{\te{ఆటో}}, \emph{āṭō}.) A taxi? (\textbf{\te{టాక్సీ}}, \emph{ṭāksī}.) Petrol? (\textbf{\te{పెట్రోలు}}, \emph{peṭrōlu}.) An ox? (\textbf{\te{ఎద్దు}}, \emph{eddu}.) A calf? (\textbf{\te{దూడ}}, \emph{dūḍa}.) A sheep? (\textbf{\te{గొర్రె}}, \emph{gorre}.) A wolf? (\textbf{\te{తోడేలు}}, \emph{tōḍēlu}.) A crocodile? (\textbf{\te{మొసలి}}, \emph{mosali}.) A thorn? (\textbf{\te{ముల్లు}}, \emph{mullu}.) A creeper? (\textbf{\te{తీగ}}, \emph{tīga}.) Bamboo? (\textbf{\te{వెదురు}}, \emph{veduru}.) Sugarcane? (\textbf{\te{చెరకు}}, \emph{ceraku}.) Worship? (\textbf{\te{పూజ}}, \emph{pūja}.) Nature? (\textbf{\te{ప్రకృతి}}, \emph{prakṛti}.) A sport? (\textbf{\te{క్రీడ}}, \emph{krīḍa}.) A player? (\textbf{\te{ఆటగాడు}}, \emph{āṭagāḍu}.) Cricket? (\textbf{\te{క్రికెట్}}, \emph{krikeṭ}.) Chess? (\textbf{\te{చదరంగం}}, \emph{cadaraṅgaṁ}.)
\end{tcolorbox}
